\begin{helper}
If the event that the size bound holds but some ordered blue pair violates
the projected-intersection estimate is negligible, then the blue implication
event has high probability.
\end{helper}
\begin{proof}
Fix \(n\) and abbreviate the two-bite probability functional by
\[
\mathbb P_n(E)=
\noderef{TwoBiteEventProbability}\bigl(n,m(n),
\noderef{TwoBiteEdgeProbability}(\beta,n),E\bigr).
\]
Let \(G_n\) be the event in which the size bound implies the blue
intersection bound for every ordered blue pair.  Let \(B_n\) be the event
that the size bound holds and at least one ordered blue pair violates the
blue bound.  The event \(G_n\) is the complement of \(B_n\).  Indeed, if
\(G_n\) fails, then the size bound holds and the universal blue estimate
fails.  Negating the universal quantifier produces an ordered pair \(b,c\)
for which the hypothesis \(b\ne c\) holds and the displayed inequality
fails.  Since the two sides of that inequality are real numbers, failure of
\(\le\) is the strict reverse inequality, so this pair witnesses \(B_n\).
Conversely, any witness to \(B_n\) makes the implication defining \(G_n\)
false.

The weights defining \(\mathbb P_n\) have exact total mass \(1\) for all
sufficiently large \(n\), by
\(\noderef{TwoBiteNaturalTotalMassOneEventually}\).  That node proves the
needed injection support \(n\le m(n)^2\) from
\(m(n)=\noderef{TwoBiteNaturalReducedVertexCount}(n)\), uses the positive
cardinality of
\(\operatorname{Fin}n\hookrightarrow(\operatorname{Fin}m(n)\times
\operatorname{Fin}m(n))\) to put \(\noderef{UniformInjectionWeight}\) in its
reciprocal branch, and then factors the whole-sample mass into two
\(\noderef{GnpGraphWeightSumOne}\) graph factors and the uniform-injection
factor.  Therefore, eventually, the finite indicator sums for complementary
events \(G_n\) and \(B_n\) add to the total mass \(1\), giving
\[
\mathbb P_n(G_n)=1-\mathbb P_n(B_n).
\]

The hypothesis says precisely that \(\mathbb P_n(B_n)\to0\), in the notation
of \(\noderef{NegligibleProbability}\).  Therefore
\[
\mathbb P_n(G_n)=1-\mathbb P_n(B_n)\to1.
\]
Unfolding \(\noderef{WithHighProbability}\), this is the desired
high-probability statement for the blue implication event.
\end{proof}
